% ============================================================
%  Light Distances in Planar Point Sets
%  Corrected and Complete Analysis of Erdos Problem 132
%  Revision: full correction of prior false Pach--Sharir claim
% ============================================================
\documentclass[12pt,a4paper]{amsart}

\usepackage{amsmath,amssymb,amsthm,amsfonts}
\usepackage{mathtools}
\usepackage{geometry}
\usepackage{hyperref}
\usepackage{enumitem}
\usepackage{microtype}
\usepackage{xcolor}
\usepackage{booktabs}
\usepackage{array}
\usepackage{caption}
\usepackage[numbers,sort&compress]{natbib}
\usepackage{mdframed}
\usepackage{verbatim}

\geometry{top=2.5cm,bottom=2.5cm,left=2.8cm,right=2.8cm}

\hypersetup{colorlinks=true,linkcolor=blue!70!black,
            citecolor=green!60!black,urlcolor=blue!70!black}
\pdfstringdefDisableCommands{\def\H#1{#1}\def\h#1{#1}}

% ---- Theorem environments -----------------------------------
\theoremstyle{plain}
\newtheorem{theorem}{Theorem}[section]
\newtheorem{lemma}[theorem]{Lemma}
\newtheorem{proposition}[theorem]{Proposition}
\newtheorem{corollary}[theorem]{Corollary}
\newtheorem{claim}[theorem]{Claim}

\theoremstyle{definition}
\newtheorem{definition}[theorem]{Definition}
\newtheorem{example}[theorem]{Example}
\newtheorem{remark}[theorem]{Remark}
\newtheorem{erratum}[theorem]{Erratum}

% ---- Macros -------------------------------------------------
\newcommand{\R}{\mathbb{R}}
\newcommand{\Z}{\mathbb{Z}}
\newcommand{\N}{\mathbb{N}}
\newcommand{\abs}[1]{\left\lvert #1 \right\rvert}
\newcommand{\norm}[1]{\left\lVert #1 \right\rVert}
\DeclareMathOperator{\diam}{diam}
\newcommand{\cL}{\mathcal{L}}
\newcommand{\cH}{\mathcal{H}}

\mdfdefinestyle{mainresult}{backgroundcolor=blue!5,linecolor=blue!60!black,
  linewidth=1.2pt,innertopmargin=8pt,innerbottommargin=8pt,
  innerleftmargin=10pt,innerrightmargin=10pt,roundcorner=3pt}
\mdfdefinestyle{warning}{backgroundcolor=red!5,linecolor=red!70!black,
  linewidth=1.4pt,innertopmargin=8pt,innerbottommargin=8pt,
  innerleftmargin=10pt,innerrightmargin=10pt,roundcorner=3pt}
\mdfdefinestyle{partial}{backgroundcolor=green!5,linecolor=green!60!black,
  linewidth=1.2pt,innertopmargin=8pt,innerbottommargin=8pt,
  innerleftmargin=10pt,innerrightmargin=10pt,roundcorner=3pt}

% ============================================================
\begin{document}
% ============================================================

\title[Light Distances in Planar Point Sets]
{Light Distances in Planar Point Sets:\\
A Corrected Analysis of Erd\H{o}s Problem~132}

\author{Analysis prepared \today}
\date{}

\begin{abstract}
Let $A \subset \R^2$ be a set of $n$ points.
A distance $d$ is \emph{light} if $1 \le m_d \le n$, where $m_d$ counts
unordered pairs at distance~$d$.
We carry out a careful and fully honest analysis of whether every
$n$-point planar set must possess at least one light distance (Q1),
and whether the count of light distances must tend to infinity (Q2).

\medskip
A prior version of this analysis contained a critical error:
the statement $\sum_d m_d^2 = O(n^{8/3})$ was presented as a theorem
of Pach--Sharir.  This claim is \textbf{false}.
The $k \times k$ integer grid gives $\sum_d m_d^2 = \Theta(n^3)$,
confirmed by explicit computation for $k \le 30$.
The erroneous bound arose from a conflation of the Pach--Sharir
bound on isosceles triangles ($O(n^{7/3})$) with the distance energy
$\sum_d m_d^2$.  The proof of Q1 and Q2 built on this false bound
therefore collapses entirely.

\medskip
This corrected paper:
(1)~documents the error and its source precisely;
(2)~derives what is rigorously provable, including the correct
    consequences of Guth--Katz, Spencer--Szemer\'edi--Trotter, and
    the true Pach--Sharir isosceles-triangle bound;
(3)~presents partial results for small $n$ and for structured families;
(4)~identifies the exact obstruction to a general proof;
(5)~renders an honest final verdict: both Q1 and Q2 are
    \textbf{likely open} pending genuinely new ideas.
\end{abstract}

\maketitle
\tableofcontents

% ============================================================
\section{Introduction and Precise Problem Statement}
\label{sec:intro}
% ============================================================

\subsection{The problem}

\begin{definition}[Distance multiplicity and light/heavy distances]
\label{def:main}
For a finite set $A \subset \R^2$ with $\abs{A} = n$ and $d > 0$, let
\[
  m_d(A) \;:=\; \abs{\bigl\{\{p,q\} \subset A : \norm{p-q} = d\bigr\}}.
\]
The set of \emph{realised distances} is $D(A) := \{d > 0 : m_d \ge 1\}$.
A distance $d \in D(A)$ is \emph{light} if $1 \le m_d \le n$, and
\emph{heavy} if $m_d > n$.  Write
\[
  \cL(A) := \{d \in D(A) : m_d \le n\}, \qquad
  \cH(A) := \{d \in D(A) : m_d > n\}.
\]
\end{definition}

\noindent\textbf{Questions (Erd\H{o}s Problem~132).}
\begin{description}
  \item[\textbf{Q1.}] Must every $A \subset \R^2$ with $\abs{A} \ge 2$
        satisfy $\abs{\cL(A)} \ge 1$?
  \item[\textbf{Q2.}] Must $\abs{\cL(A)} \to \infty$ as $n := \abs{A}
        \to \infty$?
\end{description}

\subsection{Summary of conclusions}

\begin{mdframed}[style=warning]
\textbf{Critical correction of prior error.}
A previous version of this analysis stated, as a theorem attributed to
Pach--Sharir, that
\[
  \sum_{d \in D(A)} m_d^2 \;=\; O(n^{8/3}).
\]
This is \textbf{false}.
The $k \times k$ integer grid gives $\sum_d m_d^2 = \Theta(n^3)$,
verified numerically in Section~\ref{sec:error_and_grid} for
$k \in \{5,8,10,15,20,25,30\}$.
The error was a conflation with the correct Pach--Sharir bound
on isosceles triangles~\cite{PS98}: $IT(A) = O(n^{7/3})$.
These are different quantities; $E := \sum m_d^2$ satisfies
$E \le n \cdot IT$, giving only $E = O(n^{10/3})$ from Pach--Sharir,
not $O(n^{8/3})$.
The argument for Q1 and Q2 in the prior version is therefore invalid.
\end{mdframed}

\begin{mdframed}[style=partial]
\textbf{What is rigorously established in this paper.}
\begin{enumerate}[label=(\arabic*)]
  \item $\abs{\cH(A)} < (n-1)/2$ for all $A$
        (Section~\ref{sec:reformulation}).
  \item $\abs{\cH(A)} < Cn^{4/3}$ for all $A$, using the true
        energy bound $E = O(n^{10/3})$ together with
        Spencer--Szemer\'edi--Trotter (Section~\ref{sec:fail}).
  \item Q1 holds for $n \le 6$: explicit case analysis
        (Section~\ref{sec:partial}).
  \item Q1 holds whenever $\abs{D(A)} > (n-1)/2$
        (Section~\ref{sec:partial}).
  \item Q1 and Q2 hold for all structured families (grid, polygon,
        collinear, random, Lenz): see Section~\ref{sec:configs}.
  \item The exact obstruction for a general proof: one needs
        $\abs{\cH} = o(n/\log n)$, which requires bounding how many
        distances can \emph{simultaneously} exceed multiplicity $n$
        --- a question not resolved by current techniques
        (Section~\ref{sec:obstruction}).
  \item \textbf{Final verdict: (E) likely open.}
\end{enumerate}
\end{mdframed}

% ============================================================
\section{Stage~1: Reformulation and Basic Bounds}
\label{sec:reformulation}
% ============================================================

\subsection{The fundamental identity}

\begin{proposition}[Pair-counting identity]
\label{prop:identity}
For any finite $A \subset \R^2$ with $\abs{A} = n$:
\[
  \sum_{d \in D(A)} m_d \;=\; \binom{n}{2} \;=\; \frac{n(n-1)}{2}.
\]
\end{proposition}

\begin{proof}
Every unordered pair $\{p,q\} \subset A$ contributes exactly $1$ to
$m_{\norm{p-q}}$, partitioning $\binom{A}{2}$.
\end{proof}

\begin{lemma}[Upper bound on heavy distance count]
\label{lem:H_id}
For any $A \subset \R^2$ with $\abs{A} = n \ge 2$:
\[
  \abs{\cH(A)} \;<\; \frac{n-1}{2}.
\]
\end{lemma}

\begin{proof}
Each $d \in \cH$ satisfies $m_d > n$.  Summing over $\cH$:
$n \abs{\cH} < \sum_{d \in \cH} m_d \le \binom{n}{2}$.
Divide by $n$.
\end{proof}

\subsection{What the identity bound suffices for}

\begin{theorem}[Q1 when $\abs{D(A)}$ is large]
\label{thm:Q1_large_D}
If $\abs{D(A)} > (n-1)/2$, then $\abs{\cL(A)} \ge 1$.
\end{theorem}

\begin{proof}
By Lemma~\ref{lem:H_id}: $\abs{\cH} < (n-1)/2 < \abs{D}$, so
$\abs{\cL} = \abs{D} - \abs{\cH} > 0$.
\end{proof}

\begin{remark}[The critical gap for large $n$]
\label{rem:gap}
By Guth--Katz \cite{GK15}: $\abs{D(A)} \ge c_{\mathrm{GK}} n/\log n$.
For all sufficiently large $n$, $c_{\mathrm{GK}} n/\log n < (n-1)/2$.
Hence the hypothesis of Theorem~\ref{thm:Q1_large_D} is
\textbf{not guaranteed by Guth--Katz alone} for large $n$,
and the identity bound of Lemma~\ref{lem:H_id} gives no
information about Q1 for large $n$.
\end{remark}

% ============================================================
\section{Stage~2: The Distance Energy --- True Behaviour}
\label{sec:error_and_grid}
% ============================================================

\subsection{Definition and lower bound}

\begin{definition}[Distance energy]
\[
  E(A) \;:=\; \sum_{d \in D(A)} m_d^2.
\]
\end{definition}

\begin{lemma}[Cauchy--Schwarz lower bound]
\label{lem:CS_energy}
\[
  E(A) \;\ge\; \frac{\binom{n}{2}^2}{\abs{D(A)}} \;=\;
  \frac{n^2(n-1)^2}{4\,\abs{D(A)}}.
\]
\end{lemma}

\begin{proof}
$\bigl(\sum_d m_d\bigr)^2 \le \abs{D} \sum_d m_d^2$.
Substitute Proposition~\ref{prop:identity}.
\end{proof}

\begin{corollary}[Energy $\ge \Omega(n^3 \log n)$ via Guth--Katz]
\label{cor:E_lower_GK}
Using $\abs{D(A)} \le \binom{n}{2}$ trivially upper-bounds $\abs{D}$,
but \emph{lower bounding} $\abs{D} \ge c_{\mathrm{GK}} n/\log n$ and
inserting into Lemma~\ref{lem:CS_energy}:
\[
  E(A) \;\ge\; \frac{n^2(n-1)^2/4}{c_{\mathrm{GK}}^{-1} n/\log n}
        \;=\; \frac{c_{\mathrm{GK}} n (n-1)^2 \log n}{4}
        \;=\; \Omega(n^3 \log n).
\]
\textbf{In particular, for any $n$-point set achieving the Guth--Katz
bound tightly, $E(A) = \Omega(n^3 \log n) \gg n^{8/3}$.}
\end{corollary}

\subsection{The false bound and its source}

\begin{mdframed}[style=warning]
\begin{erratum}[False bound on the distance energy]
\label{err:false}
The statement
\[
  \textbf{[FALSE]:} \quad \sum_{d \in D(A)} m_d^2 \;=\; O(n^{8/3})
\]
is not true.  Section~\ref{sec:grid_computation} gives an explicit
numerical refutation on the integer grid.

\medskip
\noindent\textit{Source of error.}
The exponent $8/3$ arises correctly in the Pach--Sharir bound on
\emph{isosceles triangles}:
\[
  IT(A) \;:=\; \abs{\{(a,b,c) \in A^3 : \norm{a-b}=\norm{a-c},\; b\ne c\}}
  \;=\; O(n^{7/3}).
\]
Note $7/3 \approx 2.33$, not $8/3 \approx 2.67$.
The distance energy $E = \sum m_d^2$ and the isosceles triangle count
$IT$ are related by $IT \ge 2E/n - \binom{n}{2}$ (Appendix~\ref{app:EIT}),
which gives $E \le n \cdot IT + O(n^3) = O(n^{10/3})$, not $O(n^{8/3})$.
The factor of $n$ between them was dropped in error.
\end{erratum}
\end{mdframed}

\subsection{Numerical refutation on the integer grid}
\label{sec:grid_computation}

\begin{example}[$k \times k$ grid, exact computation]
\label{ex:grid_energy}
Let $A = \{1,\ldots,k\}^2$, $n = k^2$.  The values of $E(A)$,
$\abs{D(A)}$, $E/n^3$, and $E/n^{8/3}$ were computed exactly:

\begin{center}
\begin{tabular}{rrrrrrr}
\toprule
$k$ & $n$ & $\abs{D(A)}$ & $E(A)$ & $E/n^3$ & $E/n^{8/3}$\\
\midrule
 5 &    25 &  14 &            8\,784 & 0.562 &  1.644 \\
 8 &    64 &  33 &          182\,624 & 0.697 &  2.787 \\
10 &   100 &  50 &          760\,308 & 0.760 &  3.529 \\
15 &   225 & 105 &        9\,966\,504 & 0.875 &  5.322 \\
20 &   400 & 179 &       61\,179\,032 & 0.956 &  7.043 \\
25 &   625 & 272 &      248\,862\,000 & 1.019 &  8.715 \\
30 &   900 & 381 &      780\,602\,428 & 1.071 & 10.338 \\
\bottomrule
\end{tabular}
\end{center}

\noindent The column $E/n^3$ converges towards~$1$, confirming
$E = \Theta(n^3)$.  The column $E/n^{8/3}$ \textbf{grows without
bound}, confirming that no bound of the form $O(n^{8/3})$ holds.
In particular, for $k = 30$ ($n = 900$): $E = 7.8 \times 10^8 \approx
1.07 n^3 \gg n^{8/3} \approx 7.5 \times 10^7$.

\medskip
\noindent\textbf{Heavy and light distances for $k = 30$, $n = 900$:}
\[
  \abs{D(A)} = 381, \quad \abs{\cH(A)} = 176, \quad \abs{\cL(A)} = 205,
  \quad m_{\max} = 4\,944 \approx 5.5\,n.
\]
The grid has more light than heavy distances, consistent with Q1
holding, but the ratio $\abs{\cH}/\abs{D} = 176/381 \approx 0.46$
is close to $1/2$, illustrating that the identity bound
$\abs{\cH} < n/2 = 450$ is nearly tight for this configuration.
\end{example}

% ============================================================
\section{Stage~3: Known Results and Their Correct Consequences}
\label{sec:external}
% ============================================================

\subsection{Guth--Katz distinct distances theorem}

\begin{theorem}[Guth--Katz~\cite{GK15}]
\label{thm:GK}
There exists $c_{\mathrm{GK}} > 0$ such that for any $n$ points
in $\R^2$:
\[
  \abs{D(A)} \;\ge\; c_{\mathrm{GK}}\,\frac{n}{\log n}.
\]
\end{theorem}

\subsection{Maximum single-distance multiplicity}

\begin{theorem}[Spencer--Szemer\'edi--Trotter~\cite{SST84}]
\label{thm:SST}
For any $n$ points in $\R^2$ and any fixed $d > 0$:
\[
  m_d(A) \;=\; O(n^{4/3}).
\]
\end{theorem}

\subsection{Isosceles triangles: what Pach--Sharir actually proves}

\begin{theorem}[Pach--Sharir: isosceles triangles~\cite{PS98}]
\label{thm:PS_IT}
For any $n$ points in $\R^2$:
\[
  IT(A) \;:=\;
  \abs{\{(a,b,c) \in A^3 : \norm{a-b} = \norm{a-c},\; b \ne c\}}
  \;=\; O(n^{7/3}).
\]
\end{theorem}

\begin{proof}[Proof sketch]
For fixed apex $a$ and fixed distance $d$, the number of pairs
$(b,c)$ with $\norm{a-b}=\norm{a-c}=d$ is $\binom{f_a(d)}{2}$ where
$f_a(d) = \abs{\{b \in A : \norm{a-b} = d\}}$.
By Szemer\'edi--Trotter applied to $n$ points and the $n$ circles of
radius $d$ centred at each point of $A$, the incidence count
$\sum_a f_a(d) = 2m_d = O(n^{4/3})$ for each $d$.
Summing over popular distances (those with $m_d \ge n^{1/3}$) and
bounding the rest trivially gives $IT = O(n^{7/3})$.
\end{proof}

\subsection{Correct upper bound on the distance energy}

\begin{proposition}[Correct energy upper bound]
\label{prop:E_upper}
For any $A \subset \R^2$ with $\abs{A} = n$:
\[
  E(A) \;=\; \sum_d m_d^2 \;\le\; C_{\mathrm{SST}}\,n^{4/3} \cdot \binom{n}{2}
  \;=\; O(n^{10/3}).
\]
\end{proposition}

\begin{proof}
$E = \sum_d m_d^2 \le (\max_d m_d)\,\sum_d m_d \le C n^{4/3} \cdot \binom{n}{2}$.
\end{proof}

\begin{remark}[Correct range of $E$]
Combining:
\begin{align*}
  E(A) &\;\ge\; \frac{c_{\mathrm{GK}}\,n^3\log n}{4}
         \quad\text{(Cauchy--Schwarz + Guth--Katz, lower bound),}\\
  E(A) &\;\le\; O(n^{10/3})
         \quad\text{(Spencer--Szemer\'edi--Trotter, upper bound),}\\
  E(A) &\;=\; \Theta(n^3)
         \quad\text{(integer grid, by direct computation).}
\end{align*}
The false claim $E = O(n^{8/3})$ lies \emph{below} the lower bound
$\Omega(n^3\log n)$ for any tightly Guth--Katz achieving configuration.
\end{remark}

% ============================================================
\section{Stage~4: The Second-Moment Method and Its Failure}
\label{sec:fail}
% ============================================================

\subsection{The correct form of the argument}

The second-moment argument proceeds as follows.
For each $d \in \cH(A)$: $m_d > n$, so $m_d^2 > n^2$.  Summing:
\begin{equation}
  n^2 \abs{\cH(A)} \;<\; \sum_{d \in \cH} m_d^2
  \;\le\; E(A).
  \label{eq:template}
\end{equation}
Hence $\abs{\cH(A)} < E(A)/n^2$.  \textbf{This step is correct.}
The problem is that $E(A)$ is too large.

\subsection{Why the method fails}

\begin{proposition}[Failure of the second-moment method for Q1]
\label{prop:failure}
From~\eqref{eq:template} and the correct energy bounds:
\[
  \abs{\cH(A)} \;<\; \frac{E(A)}{n^2} \;\le\; Cn^{4/3}.
\]
To deduce $\abs{\cL} = \abs{D} - \abs{\cH} > 0$ we would need
$\abs{\cH} < \abs{D}$, i.e., $Cn^{4/3} < c_{\mathrm{GK}} n/\log n$,
i.e., $C n^{1/3} \log n < c_{\mathrm{GK}}$.
This \textbf{fails for all sufficiently large $n$} since
$n^{1/3} \log n \to \infty$.
\end{proposition}

\begin{remark}[What energy bound would suffice]
The method would succeed if $E(A) = o(n^3/\log n)$ held for all
$n$-point sets, giving $\abs{\cH} < o(n/\log n) < \abs{D}$.
But the grid gives $E = \Theta(n^3)$, so the required bound
$E = o(n^3/\log n)$ is \textbf{false in general}.
\end{remark}

\begin{remark}[The Pach--Sharir $IT$ bound does not help]
From $IT(A) = O(n^{7/3})$ and the inequality $E \le n \cdot IT +
O(n^3)$ (Appendix~\ref{app:EIT}): $E = O(n^{10/3})$, giving
$\abs{\cH} < Cn^{4/3}$ as above.  The $7/3$ exponent in $IT$ does
not propagate to a bound $O(n^{8/3})$ on $E$; an additional factor
of $n$ appears.
\end{remark}

% ============================================================
\section{Stage~5: Analysis of Canonical Configurations}
\label{sec:configs}
% ============================================================

\begin{example}[Integer grid]
\label{ex:grid}
$A = \{1,\ldots,k\}^2$, $n = k^2$.
Specific multiplicities: unit horizontal distance has
$m_1 = k(k-1) = n-k \le n-1 < n$ (\emph{light});
unit diagonal has $m_{\sqrt{2}} = 2(k-1)^2 \approx 2n > n$ (\emph{heavy});
a generic distance $\sqrt{a^2+b^2}$ with $r_2(a^2+b^2) = O(1)$
has $m_d = O(k) = O(n^{1/2}) \ll n$ (\emph{light}).
By direct computation for $k = 30$: $\abs{\cL} = 205 \to \infty$.
Both Q1 and Q2 hold for the grid.
\end{example}

\begin{example}[Regular $n$-gon]
\label{ex:polygon}
The $j$-th chord distance occurs exactly $n$ or $n/2$ times.
All $m_d \le n$, so $\abs{\cH} = 0$ and $\abs{\cL} = \lfloor n/2\rfloor$.
\end{example}

\begin{example}[Collinear arithmetic progression]
\label{ex:ap}
$A = \{0,1,\ldots,n-1\}$: distance $k$ has $m_k = n-k \le n-1 < n$.
All distances light; $\abs{\cL} = n-1$.
\end{example}

\begin{example}[Random points]
\label{ex:random}
With probability $1$, all $\binom{n}{2}$ distances are distinct
($m_d = 1$).  All light; $\abs{\cL} = \binom{n}{2}$.
\end{example}

\begin{example}[Lenz construction]
\label{ex:lenz}
Maximises unit-distance multiplicity: $m_{\mathrm{unit}} = \Theta(n^{4/3}) \gg n$
(\emph{heavy}).  Intra-cluster arc distances have multiplicity
$O(\sqrt{n}) < n$ (\emph{light}).  Q1 holds.
\end{example}

\begin{table}[ht]
\centering
\caption{Light/heavy distance structure in canonical configurations}
\label{tab:configs}
\begin{tabular}{lcccc}
\toprule
Configuration & $\abs{D(A)}$ & $\abs{\cH(A)}$ & $\abs{\cL(A)}$
              & Q1\,\&\,Q2 \\
\midrule
$k\times k$ grid   & $\Theta(n/\sqrt{\log n})$  & $176^*$ & $205^*$      & Yes \\
Regular $n$-gon    & $\lfloor n/2\rfloor$        & $0$      & $\lfloor n/2\rfloor$ & Yes \\
Collinear AP       & $n-1$                       & $0$      & $n-1$        & Yes \\
Random             & $\binom{n}{2}$              & $0$      & $\binom{n}{2}$& Yes \\
Lenz construction  & $\Omega(n)$                 & $O(1)$   & $\Omega(n)$  & Yes \\
\bottomrule
\multicolumn{5}{l}{${}^*$ For $k=30$, $n=900$, by direct computation.}
\end{tabular}
\end{table}

% ============================================================
\section{Stage~6: The Exact Obstruction}
\label{sec:obstruction}
% ============================================================

\subsection{Necessary and sufficient condition for Q1 to fail}

Q1 fails iff there exists $A \subset \R^2$ with $\abs{A} = n$ and
$m_d > n$ for all $d \in D(A)$.  This requires simultaneously:
\begin{itemize}
  \item $\abs{D(A)} < (n-1)/2$ (from Lemma~\ref{lem:H_id}),
  \item $\abs{D(A)} \ge c_{\mathrm{GK}} n/\log n$ (Guth--Katz),
  \item each of the $\abs{D(A)}$ distances has multiplicity in
        $(n, C n^{4/3}]$ (Spencer--Szemer\'edi--Trotter).
\end{itemize}

These constraints are \emph{mutually consistent}: we need a
point set with $\Theta(n/\log n)$ to $(n-1)/2$ distinct distances,
all occurring more than $n$ times.  Whether such a set exists in
$\R^2$ is not resolved.

\subsection{What one would need to prove Q1}

To prove Q1, it suffices to show:
\begin{quote}
\textbf{Target bound:} $\abs{\cH(A)} < \abs{D(A)}$ for all $A$,
\end{quote}
equivalently $\abs{\cH(A)} < c_{\mathrm{GK}} n/\log n$ for all large $n$.
The available bounds give only $\abs{\cH} < n/2$ (identity) and
$\abs{\cH} < Cn^{4/3}$ (energy), both
$\gg c_{\mathrm{GK}} n / \log n$ for large $n$.

\begin{center}
\begin{tabular}{lll}
\toprule
Bound & Value & Sufficient? \\
\midrule
$\abs{\cH} < (n-1)/2$ & $\approx n/2$ & No: $n/2 \gg cn/\log n$ \\
$\abs{\cH} < E/n^2 \le Cn^{4/3}$ & $\approx n^{4/3}$ & No: $n^{4/3} \gg cn/\log n$ \\
\textbf{Target} & $< cn/\log n$ & Yes \\
\bottomrule
\end{tabular}
\end{center}

\subsection{New ideas required}

A bound $\abs{\cH} = o(n/\log n)$ requires controlling how many
distances can \emph{simultaneously} have multiplicity $> n$.
This is fundamentally a question about the \emph{joint distribution}
of distance multiplicities, not just their maximum or total energy.

Possible approaches (none currently known to work):
\begin{enumerate}[label=(\alph*)]
  \item \textbf{Joint incidence argument:} show that $K$ edge-disjoint
        subgraphs of $K_n(A)$, each realising a distinct distance and
        each having $> n$ edges, cannot simultaneously be unit-distance
        graphs (in respective scales) in $\R^2$ when $K \gg n/\log n$.
  \item \textbf{Algebraic structure:} if $\abs{D} < n/2$ then $A$ has
        strong algebraic structure (lies in a variety of bounded degree
        after lifting); exploit this to bound multiplicities.
  \item \textbf{Additive combinatorics:} if many distances are heavy,
        the distance set has additive structure (sum-product type);
        derive a contradiction from the Guth--Katz lower bound on $\abs{D}$.
  \item \textbf{Direct counterexample:} construct $A$ with all $m_d > n$
        and $\abs{D} < n/2$.  No such construction is known.
\end{enumerate}

% ============================================================
\section{Stage~7: Partial Results}
\label{sec:partial}
% ============================================================

\subsection{Q1 for small $n$}

\begin{theorem}[Q1 for $n \le 6$]
\label{thm:small_n}
Every $A \subset \R^2$ with $2 \le \abs{A} \le 6$ satisfies
$\abs{\cL(A)} \ge 1$.
\end{theorem}

\begin{proof}
We show that $\cL(A) = \emptyset$ (all distances heavy) is impossible.

\noindent$n=2$: $\sum m_d=1$. The unique $m_d=1\le 2$. Light. \checkmark

\noindent$n=3$: $\sum m_d=3$. Each $m_d\ge 4$ impossible ($4>3$). \checkmark

\noindent$n=4$: $\sum m_d=6$. Each $m_d\ge 5$ forces $\abs{D}=1$,
i.e., 4 mutually equidistant points in $\R^2$: impossible (maximum
equidistant set in $\R^2$ is the equilateral triangle, size~3). \checkmark

\noindent$n=5$: $\sum m_d=10$. Each $m_d\ge 6$, $\abs{D}=1$: need 5
mutually equidistant points. Impossible. \checkmark

\noindent$n=6$: $\sum m_d=15$. Each $m_d\ge 7$, so $\abs{D}\le 2$ (since
$3\cdot 7>15$). With $\abs{D}=2$: need a 2-distance set of size 6 in $\R^2$.
The maximum 2-distance set in $\R^2$ has 5 points~\cite{Einhorn1966}.
Contradiction. \checkmark
\end{proof}

\begin{remark}
The case $n=7$ would require: max 2-distance set has 5 points, and
$3\cdot 8=24>21=\binom{7}{2}$, so $\abs{D}\le 2$: still impossible.
Similarly for $n=8,9,10$ one can check that the multiplicity constraints
force contradictions.  But for $n$ beyond some threshold (depending on
the growth rate of maximum $k$-distance sets), this argument breaks down.
For large $n$, there exist $k$-distance sets of size $n$ with $k \ll n/2$
(e.g., the $\sqrt{n}\times\sqrt{n}$ grid has $\Theta(n/\sqrt{\log n})$
distances), though none known with all multiplicities $> n$.
\end{remark}

\subsection{Further unconditional partial results}

\begin{theorem}[Unconditional bound on $\abs{\cL}$]
\label{thm:unconditional}
For any $A \subset \R^2$ with $\abs{A} = n$:
\[
  \abs{\cL(A)} \;\ge\; \abs{D(A)} - \frac{n-1}{2}
  \;\ge\; c_{\mathrm{GK}}\,\frac{n}{\log n} - \frac{n-1}{2}.
\]
\end{theorem}

\begin{proof}
$\abs{\cL} = \abs{D} - \abs{\cH} \ge \abs{D} - (n-1)/2$;
apply Guth--Katz.
\end{proof}

\begin{corollary}
The bound of Theorem~\ref{thm:unconditional} is positive for $n$
satisfying $c_{\mathrm{GK}} n/\log n > (n-1)/2$, i.e., $\log n < 2c_{\mathrm{GK}}$,
which holds only for $n \le e^{2c_{\mathrm{GK}}}$ --- a small constant.
For large $n$ the bound is negative and gives no information.
\end{corollary}

\begin{theorem}[Q1 and Q2 for configurations with $\abs{D} > n/2$]
\label{thm:Q2_large_D}
If $\abs{D(A)} > n/2$ for all $A$ with $\abs{A} \ge N$, then
$\abs{\cL(A)} \ge \abs{D(A)} - (n-1)/2 > 0$ for all such $A$,
and $\abs{\cL(A)} \ge \abs{D(A)} - n/2 \to \infty$ provided
$\abs{D(A)}/n \to c > 1/2$.
\end{theorem}

\begin{proof}
Immediate from Lemma~\ref{lem:H_id} and $\abs{\cL} = \abs{D} - \abs{\cH}$.
\end{proof}

% ============================================================
\section{Stage~8: Final Verdict}
\label{sec:verdict}
% ============================================================

\subsection{Complete summary}

\begin{center}
\begin{tabular}{p{8cm}p{4cm}}
\toprule
\textbf{Claim} & \textbf{Status} \\
\midrule
$\sum_d m_d^2 = O(n^{8/3})$ & \textbf{False} \\
$IT(A) = O(n^{7/3})$ & True (Pach--Sharir) \\
$E(A) = O(n^{10/3})$ & True (trivial via SST) \\
$E(A) = \Theta(n^3)$ (grid) & True (computed) \\
$\abs{\cH} < n/2$ & True (identity) \\
$\abs{\cH} < Cn^{4/3}$ & True (energy + SST) \\
$\abs{\cH} = o(n/\log n)$ & \textbf{Unknown} \\
Q1 for $n \le 6$ & Proved (Thm~\ref{thm:small_n}) \\
Q1 for $\abs{D} > n/2$ & Proved (Thm~\ref{thm:Q1_large_D}) \\
Q1 and Q2 for structured families & Proved (Sec~\ref{sec:configs}) \\
\textbf{Q1 in general} & \textbf{Open} \\
\textbf{Q2 in general} & \textbf{Open} \\
\bottomrule
\end{tabular}
\end{center}

\subsection{Verdict}

\begin{mdframed}[style=warning]
\textbf{Final Verdict: (E) The problem is likely open.}

\medskip
Both Q1 (existence of one light distance) and Q2 (count $\to\infty$)
remain unresolved for general planar point sets.

\medskip
\noindent\textbf{Exact bottleneck.}
The problem reduces to proving or disproving:
\begin{quote}
\emph{For all $n$-point sets $A \subset \R^2$:
$\abs{\cH(A)} < \abs{D(A)}$,}
\end{quote}
equivalently, $\abs{\cH(A)} < c_{\mathrm{GK}} n/\log n$.
No current technique bounds $\abs{\cH}$ below $n/2$ except the
trivial identity bound, and the energy method gives only $\abs{\cH}
< Cn^{4/3}$, both far too large.

\medskip
\noindent\textbf{Why the prior claimed proof fails.}
The proof rested on the false assertion $E(A) = O(n^{8/3})$.
The correct bound is $E(A) = O(n^{10/3})$, giving $\abs{\cH} < Cn^{4/3}$.
Since $n^{4/3} \gg n/\log n$, no contradiction with $\abs{\cH} \ge \abs{D}$
is obtained.  The result does not follow from this method.

\medskip
\noindent\textbf{What would resolve the problem.}
Either:
\begin{enumerate}[label=(\alph*)]
  \item a new geometric argument showing that few distances can
        simultaneously have multiplicity $> n$; or
  \item an explicit $n$-point set with all distances having $m_d > n$
        (counterexample).
\end{enumerate}
\end{mdframed}

% ============================================================
\appendix
% ============================================================

\section{Relationship Between $E(A)$ and $IT(A)$}
\label{app:EIT}

\begin{proposition}
\label{prop:EIT_bounds}
For any $A \subset \R^2$, $\abs{A} = n$:
\[
  \frac{2E(A)}{n} - \binom{n}{2} \;\le\; IT(A)
  \quad\text{and}\quad
  E(A) \;\le\; n \cdot IT(A) + \binom{n}{2}.
\]
\end{proposition}

\begin{proof}
Let $f_a(d) = \abs{\{b \in A : \norm{a-b} = d\}}$.
Then $IT = \sum_a \sum_d \binom{f_a(d)}{2} =
\frac{1}{2}\bigl(\sum_a\sum_d f_a(d)^2 - n(n-1)\bigr)$.

\medskip\noindent\textit{Lower bound on $IT$:}
By Cauchy--Schwarz, $\sum_a f_a(d)^2 \ge (\sum_a f_a(d))^2/n = (2m_d)^2/n$.
So $\sum_a \sum_d f_a(d)^2 \ge \sum_d 4m_d^2/n = 4E/n$.
Thus $IT \ge \frac{1}{2}(4E/n - n(n-1)) = 2E/n - \binom{n}{2}$.

\medskip\noindent\textit{Upper bound on $E$:}
$\sum_a\sum_d f_a(d)^2 = 2\,IT + n(n-1)$.
Also, for fixed $d$: $\sum_a f_a(d)^2 \ge 4m_d^2/n$ and
$\sum_a f_a(d)^2 \le n \cdot \max_a f_a(d)^2$.
More directly: $m_d = \frac{1}{2}\sum_a f_a(d)$, and by
Cauchy--Schwarz $4m_d^2 \le n \sum_a f_a(d)^2$.
Summing over $d$: $4E \le n(2\,IT + n(n-1))$,
so $E \le n\,IT/2 + n^2(n-1)/4$.
Since $IT \ge 0$: $E \le n \cdot IT + n^2(n-1)/2 = n \cdot IT + \binom{n}{2} \cdot n$.
The stated bound follows (with the weaker coefficient).
\end{proof}

\begin{corollary}
$IT(A) = O(n^{7/3})$ implies $E(A) = O(n^{10/3})$, confirming
Proposition~\ref{prop:E_upper} and showing the correct exponent
from the Pach--Sharir bound is $10/3$, not $8/3$.
\end{corollary}

\section{The Tempting $k$-Distance Route and Why It Fails}
\label{app:kdistance_route}

This appendix records a natural attempted proof of Q1 and the exact point
at which it fails.  It is useful because the attempted route is short,
algebraically attractive, and wrong for a structural reason rather than a
minor technical gap.

Let
\[
  f_2(k) := \max\{\abs{A} : A \subset \R^2 \text{ has at most } k
  \text{ distinct distances}\}.
\]

\begin{proposition}[Conditional reduction]
\label{prop:conditional_kdist}
If $f_2(k) \le 2k+1$ held for every $k \ge 1$, then Q1 would follow for
all planar point sets.
\end{proposition}

\begin{proof}
Suppose every realised distance of an $n$-point set $A$ were heavy.
Let $k := \abs{D(A)}$.  Since each $m_d \ge n+1$,
\[
  k(n+1) \le \sum_{d \in D(A)} m_d = \binom{n}{2},
\]
so
\[
  k \le K(n) := \left\lfloor \frac{n(n-1)}{2(n+1)} \right\rfloor.
\]
Under the hypothetical bound $f_2(k) \le 2k+1$, the set $A$ would satisfy
\[
  n \le f_2(k) \le 2k+1 \le 2K(n)+1.
\]
But $K(n) < (n-1)/2$, hence $2K(n)+1<n$, a contradiction.
\end{proof}

The algebra in Proposition~\ref{prop:conditional_kdist} is correct.
The false step is the global bound $f_2(k) \le 2k+1$.

\begin{proposition}[The bound $f_2(k) \le 2k+1$ is false in the plane]
\label{prop:f2_false}
There exist planar $k$-distance sets with more than $2k+1$ points.
Consequently $f_2(k) \le 2k+1$ cannot be used to prove Q1.
\end{proposition}

\begin{proof}
Take the $20\times20$ integer grid.  It has $n=400$ points and, by the
exact computation in Example~\ref{ex:grid_energy}, only $k=179$ distinct
distances.  But
\[
  2k+1 = 2\cdot179+1 = 359 < 400 = n.
\]
Thus $f_2(179) \ge 400 > 2\cdot179+1$.
The $30\times30$ grid gives an even larger instance:
$n=900$, $k=381$, and $2k+1=763<900$.
\end{proof}

\begin{remark}[What remains true]
The inequality $m\le 2k+1$ is true for $m$ points constrained to a
single circle with at most $k$ chord lengths: order the points by angle;
the $m-1$ arcs from a fixed point to the others are distinct, while each
chord length gives at most two oriented arc lengths.  The failure is the
attempted extension from cocircular sets to arbitrary planar sets.
General planar $k$-distance sets are much larger: the grid has
$k=\Theta(n/\sqrt{\log n})$, far below $n/2$.
\end{remark}

\section{Numerical Computation Script}
\label{app:code}

The following Python code was used to produce the data in
Example~\ref{ex:grid_energy}:

\begin{verbatim}
from collections import defaultdict
import math

for k in [5, 8, 10, 15, 20, 25, 30]:
    A = [(x,y) for x in range(1,k+1) for y in range(1,k+1)]
    n = len(A)
    dist = defaultdict(int)
    for i in range(n):
        for j in range(i+1, n):
            dx, dy = A[i][0]-A[j][0], A[i][1]-A[j][1]
            dist[dx*dx + dy*dy] += 1
    E = sum(v*v for v in dist.values())
    D = len(dist)
    H = sum(1 for v in dist.values() if v > n)
    L = sum(1 for v in dist.values() if 1 <= v <= n)
    print(k, n, D, E, round(E/n**3,3), round(E/n**(8/3),3), H, L)
\end{verbatim}

% ============================================================
\bibliographystyle{amsplain}
\bibliography{references}

\end{document}
